
\subsection{Motivation}
The pervasive integration of artificial intelligence in recent years, specifically through the deployment of Large Language Models and the rise of autonomous financial agents, has fundamentally restructured the operational paradigms of financial services. Recent research across AI application domains has also indicated that practical deployment criteria, such as computational efficiency in domain adaptation, robustness to prompt-level interaction conditions, and domain-specific evaluation design, may be at least as important as benchmark performance metrics \cite{codes2026,cai2026tone}. Although commercial investment has facilitated significant progress in automated advisory, the application of such technologies to the risk management of complex derivatives remains a critical frontier. A persistent challenge in current financial applications driven by artificial intelligence is the methodological divergence between pricing model calibration and actual hedging performance. Traditional frameworks frequently rely on static diagnostics that fail to incorporate the operational realities of market frictions and equilibrium effects in incomplete markets. This research addresses this divergence by introducing a reinforcement learning framework that shifts the hedging objective from the traditional minimization of errors to the optimization of shortfall probability.

The study utilizes agents based on neural networks within modified QLBS and novel RLOP architectures to integrate risk aversion and market frictions into the decision process. Empirical results demonstrate that these reinforcement learning agents provide a systematic cost advantage over traditional Delta hedging by establishing a decision loop sensitive to costs. In the presence of substantial transaction costs, the modified QLBS agent prioritizes stability while the RLOP model mitigates margin pressure and liquidity demand. This resilience was validated during the 2020 crash, in which RLOP systematically reduced exposure to manage extreme stress. Distributional analysis confirms that this approach achieves material improvements in downside control by reducing extreme losses after costs, ensuring a robust equilibrium during operationally salient regimes of volatility.

This paper delivers three primary research contributions: (1) This paper extends the QLBS framework by embedding shortfall probability into the reward structure, resolving the decoupling between static calibration and execution. It demonstrates that IVRMSE-based diagnostics favor parametric models but fail to reflect hedging quality under frictions, necessitating a survival-centric strategy. (2) The novel RLOP model is introduced for superior tail-risk resilience. By prioritizing hedging success frequency over loss magnitude, RLOP provides material improvements in downside control. Distributional Expected Shortfall (ES) analysis confirms it significantly reduces extreme after-cost losses during operationally salient regimes, like the 2020 crash. (3) This study establishes a bidirectional selection framework supported by cost-risk maps and full-distribution tail risk analysis. This paper proves RL policies achieve a systematic cost advantage and turnover reduction. While QLBS acts as a cost-aware stabilizer, RLOP manages margin pressure, ensuring a robust equilibrium that avoids selection bias through multi-dimensional risk-cost metrics.

The remainder of this paper is structured as follows: Section 1.2 provides a comprehensive review of the relevant literature, Section 2 shows the reinforcement learning framework, Section 3 introduces the modified QLBS model and the novel RLOP model, Section 4 and 5 present their empirical analysis under varying risk and cost scenarios, and Section 6 concludes this paper.

\subsection{Literature Review}
%Option pricing and hedging remain core challenges in quantitative finance. 
% While the Black-Scholes-Merton (BSM) framework \cite{black1973pricing} provides a foundational option valuation model using equilibrium conditions of a hedged portfolio, 
 %\cite{black1973pricing} provide a foundational option valuation model where perfect replication is achieved under frictionless markets with continuous-time trading, and subsequent research has extended BSM to adapt to more complex market realities (\cite{fan2022empirical,li2025analytic,golbabai2013superconvergence}). Nevertheless, real-world markets involve transaction costs and discrete-time trading. Classical options pricing models like Black-Scholes and stochastic volatility models are typically used for static calibration of the market’s implied volatility surface. However, as \cite{lassance2018comparison} have shown, these models do not necessarily produce better results when used to calculate hedging performance. This indicates that increasing sophistication in pricing models does not yield better hedging results as they represent two different objectives of pricing. Therefore, traditional diagnostic methods such as IVRMSE do not adequately reflect the actual hedging quality due to frictions in the real world.
 
%The evolution of option hedging reveals a persistent tension between static calibration and realized performance.%
While foundational models such as the Black-Scholes-Merton (BSM) framework (\cite{black1973pricing}), the Merton’s jump-diffusion (\cite{merton1976discontinuous}) , and the Heston’s stochastic volatility (\cite{heston1993closed}) prioritize minimizing Implied Volatility (IV) surface fitting errors, a superior static fit does not inherently translate into effective risk mitigation in the presence of transaction costs. This pricing-hedging disconnect, initially discussed in classic works like Lassance and Vrins (\cite{lassance2018comparison}), has been elevated to a primary research focus in recent years. Specifically, \cite{franccois2025difference} theoretically formalize how IV surface information can become decoupled from optimal hedging decisions under market frictions, aligning with the findings of \cite{hey2023cost} on the misspecifying price. Furthermore, recent evidence from \cite{bayer2025deep}, \cite{brugiere2025model}, and \cite{poddar2026uncertainty} confirms that minimizing the implied volatility error often fails to ensure path-dependent stability, leading to a divergence between diagnostic precision and survival outcomes. Coupled with the recent analytical perspectives of \cite{li2025analytic} and tail risk optimization and constraint handling explored by \cite{park2025robust}, these developments directly motivate our investigation into why traditional diagnostic errors poorly predict after-cost hedging quality.

Reinforcement learning has recently gained attention as a data-driven approach to hedging because it can directly optimize realized performance rather than rely on model-specific assumptions. Halperin \cite{halperin2020qlbs,halperin2019qlbs} introduces the QLBS framework, where option hedging is formulated as a discrete-time reinforcement learning problem. Buehler et al. \cite{buehler2019deep} further advanced RL-based hedging through Deep Hedging, employing neural networks to learn hedging policies under market frictions and convex risk measures. Later studies extended this paradigm to quantile-based objectives \cite{cao2023gamma}, empirical training on observed option-market data \cite{mikkila2023empirical}, time-consistent recursive expectile risk measures \cite{marzban2023deep}, convex market-impact modeling \cite{neagu2024deep}, and cross-algorithm benchmarking \cite{neagu2025deep,pickard2023deep}. Although these frameworks allow for the modeling of transaction costs, they focus primarily on replication accuracy under error-based reward functions, which can incentivize frequent trading and reduce overall cost effectiveness.
%Reinforcement learning has recently emerged as a new method of data-driven hedging, since it allows for direct optimization of realized performance instead of relying on model-specific assumptions. \cite{halperin2020qlbs,halperin2019qlbs} introduces a QLBS framework, which formulates hedging of options as a discrete-time reinforcement learning problem. \cite{buehler2019deep}further advanced RL-based hedging by introducing Deep Hedging. It is based on the use of neural networks to learn hedging policies so that transaction costs and market frictions can be included when optimizing under different risk measures (e.g., Expected Shortfall (ES)). Although these frameworks allow for the modeling of transaction costs, they still primarily focus on accuracy of replication with respect to the error-based reward functions. This can incentivize frequent trading and reduce overall cost effectiveness.

Although these reinforcement learning methods have incorporated risk measures such as Expected Shortfall \cite{buehler2019deep} and dynamic expectile risk measures \cite{marzban2023deep}, they have concentrated on the magnitude of the tail loss (expected value conditional on being in the worst $\alpha$-tail) rather than on the probability of achieving any loss at all. Föllmer and Leukert \cite{follmer2000efficient} showed that the optimal hedge under transaction costs should focus on the minimization of shortfall rather than replication error, trading off replication cost against downside protection. The 2020 COVID-19 crash underscored this distinction: operationally, strategies that reduce loss frequency are more important than those that only reduce loss magnitude. Therefore, it is necessary to introduce shortfall-aware objectives that directly optimize the probability of hedging against future shortfalls into RL strategies.
%Although these reinforcement learning method have incorporated risk measures like Expected Shortfall, they have concentrated on the magnitude of the tail loss (expected value conditional on being in the worst $\alpha$-tail) rather than on the probability of achieving any loss at all. \cite{follmer2000efficient} showed that the optimal hedge under transaction costs should focus on the minimization of shortfall rather than replication error, trading off replication cost against downside protection. The magnitude of this distinction is made clear by the experience of 2020 and the extreme volatility caused by the COVID-19 pandemic. In particular, the crisis proved the need for ``survival'' strategies. This kind of strategies focus on reducing tail risk in order to create a business that will last and ensure that the business has a level of operational stability. Therefore, it is necessary to introduce shortfall-aware objectives to directly optimize the probability of hedging against future shortfalls into RL strategies.


Moreover, existing studies report aggregate after-cost outcomes without separating execution intensity from replication dispersion. This paper introduces a cost-risk map that seperates transaction cost from pre-cost replication quality,
enabling a more precise evaluation of learned strategies.
% enabling a more precise evaluation of learned hedging strategies.